\section{Renormalization in the space of coordinate values}
\label{re:section}

We now turn the harmonic coordinates into a bi-Lipschitz chart.  The
construction takes place in the ordinary Euclidean space of their values.
In particular, the differentiations below do not involve a pointwise inverse
of the differential of the harmonic map.

\subsection{Notation and quantitative inputs}
\label{re:input-data}

Write $r_i=2^{-i}$, $P_x=u(x)$, and
\[
  \Omega=u(B_1(p)).
\]
By Lemma~\ref{co:open-coordinates}, $\Omega$ is open and $u$ is a
homeomorphism from $B_1(p)$ onto $\Omega$.  We use the positive definite
matrices $Q_i(x)$ and $H_i(x)=Q_i(x)^{-1}$ of
Lemma~\ref{co:heat-matrices}.  For a positive definite matrix $H$,
put $|\xi|_H=(\xi^T H\xi)^{1/2}$.  Vectors in the source value space
are measured with $H_i(x)$; covectors in that space are measured with
$Q_i(x)$, its dual matrix.  If $T$ is an $m$-linear map with covector
values, our operator norm is
\[
 \|T\|_{H\to Q}
 =\sup\bigl\{|T(\xi_1,\ldots,\xi_m)|_Q:
                 |\xi_a|_H\le1\ (1\le a\le m)\bigr\}.
\]
The notation $\|T\|_{H\to\R^n}$ has the analogous meaning with the
Euclidean target norm.  For $m=0$ these conventions mean the target norm
of the value.  We denote ordinary Fr\'echet derivatives in value space
by $\mathrm D^m$.  Whenever a derivative norm at $P_x$ uses $H_i(x)$ or
$Q_i(x)$, these matrices are frozen: they are not differentiated.

The fixed constants from the preceding sections, including all required
enlargements of balls, will be chosen before the final rescaling.  We
may then use a small parameter $\varepsilon>0$ tending to zero along
that sequence of rescalings.  Increasing this parameter by a fixed
constant factor if necessary, the following bounds hold:
\begin{equation}
 \label{re:monotone-metrics}
 H_i(x)\le H_{i+1}(x)\le(1+\varepsilon)H_i(x),
 \qquad Q_0(x)=\Id+o(1).
\end{equation}
The matrices $R_i(y)^{-1}$ and $H_i(x)$ are uniformly comparable if
$d(x,y)\le C r_i$, for each of the fixed constants $C$ used below.
Moreover, normalized harmonic approximation gives the separation and
central-ellipsoid conclusions of
Lemmas~\ref{co:harmonic-normalization} and
\ref{co:open-coordinates}.  In particular, there are fixed $a>0$ and
$C_{\mathrm{cmp}}\ge1$ such that, whenever the balls involved are
interior,
\begin{equation}
 \label{re:ellipsoid-input}
 \{P:|P-P_x|_{H_i(x)}<a r_i\}
       \subset u(B_{r_i}(x)),
 \qquad H_i(y)\le C_{\mathrm{cmp}}H_i(x)
       \quad(d(x,y)\le r_i).
\end{equation}
Here and below, the harmonic map is defined on the larger balls supplied
by the preceding sections, even when an image in
\eqref{re:ellipsoid-input} is not contained in $\Omega$.

For a fixed sufficiently large $L$, recall the positive semidefinite
matrices
\begin{equation}
 \label{re:patch-matrix}
 D_{i,y}=R_i(y)^{-1}
       \bigl[r_i^{2-n}M(B_{Lr_i}(y))\bigr]R_i(y)^{-1}.
\end{equation}
Lemma~\ref{co:cumulative-comparison} applies to boundedly many
centers at distance $O(r_i)$ from a testing point, with weights in
$[0,1]$ and with at least one weight equal to one whose mass ball contains
$B_{r_i}(x)$.  With the notation $Z_i$ for the unweighted sum and $S_i$
for the weighted sum, it gives a fixed constant $C_*\ge1$ such that
\begin{equation}
 \label{re:cumulative-input}
 \begin{aligned}
 Z_i(x)&\le\varepsilon H_i(x),\\
 \Id+\sum_{i<k}Z_i(x)&\le C_*H_k(x),\\
 H_k(x)&\le C_*\left(\Id+\sum_{i<k}S_i(x)\right).
 \end{aligned}
\end{equation}
Enlarging $C_*$ if needed, these statements include $k=0$.

By Proposition~\ref{pi:directional-pinching}, for these fixed choices
of $L$ and center separations there are functions $E_i$ on a fixed
domain containing $B_1(p)$ such that
\begin{equation}
 \label{re:pinching-input}
 \sup_{x\in B_1(p)}\sum_{i=0}^\infty E_i(x)\le\varepsilon,
 \qquad E_i(x)\ge e_i^*>0,
\end{equation}
where $e_i^*$ is independent of $x$ at each fixed $i$, and
\begin{equation}
 \label{re:directional-input}
 D_{i,y}[P_x-P_z]
       \le C r_i^2\bigl(E_i(x)+E_i(z)\bigr)
 \quad\text{if }d(x,y),d(z,y)\le C_0r_i.
\end{equation}
We use $D[\xi]=\xi^TD\xi$ for quadratic forms.  The positive floor in
\eqref{re:pinching-input} is supplied by the exponentially decreasing
positive term in the construction of $E_i$.  Such a term can always
be added with arbitrarily small sum.  It is used only to select fixed
centers for the patches.

All later constants depend on dimension and the fixed preceding
constants, but not on the scale, the testing point, or the number of
scales.  Constants denoted by $C_A$ may additionally depend on a
bootstrap constant $A$ and on the jet constants fixed after $A$.

\subsection{Smooth fields with directional control}

For each $i$ choose a maximal $r_i$-separated subset
$\mathcal Y_i\subset B_1(p)$.  Properness and local doubling imply
that this family is finite; maximality implies that every point of
$B_1(p)$ is at distance at most $r_i$ from some member.

Fix a smooth function $\psi:\R^n\to[0,1]$ equal to one on
$\{|X|\le2\}$ and supported in $\{|X|<3\}$, with
\begin{equation}
 \label{re:bump-jets}
 \|\mathrm D^m\psi\|_{\infty}
       \le C_b^{m+1}(m!)^s\qquad(m\ge0)
\end{equation}
for a fixed $s>1$.  One may, for example, take $s=2$.  To recall why
such bounds are available, the function $e^{-1/t}$ for $t>0$, extended
by zero for $t\le0$, has on each fixed bounded interval derivative
bounds $C^{m+1}(m!)^2$.  Cauchy's estimate on a complex disk of radius
a fixed small multiple of $t$ bounds its $m$th derivative by
$m!(C/t)^m e^{-c/t}$, and optimization in $t$ gives this assertion.
Products, integration, and composition with a fixed quadratic polynomial
give a radial cutoff with the specified radii and the same factorial
order, by the finite product and chain formulas.

Define
\[
 \phi_{i,y}(P)=
 \psi\left(r_i^{-1}R_i(y)^{-1/2}(P-P_y)\right).
\]
For $x\in B_1(p)$ let
\[
 \mathcal A_i(x)=\{y\in\mathcal Y_i:P_x\in\supp\phi_{i,y}\}.
\]
The harmonic separation estimates show that every member of
$\mathcal A_i(x)$ satisfies $d(x,y)\le C_0r_i$ for a fixed $C_0$.
Indeed, fix a growth exponent $\alpha=1/16$ in~\eqref{co:slow-R}.
For $r_i\le t=d(x,y)\le1$, choose $j\le i$ with $t\le r_j<2t$.
The separation estimate at scale $r_j$, followed by inverse matrix
comparison in~\eqref{co:slow-R}, gives
\[
 |P_x-P_y|_{R_i(y)^{-1}}
 \ge c(r_i/r_j)^\alpha |P_x-P_y|_{R_j(y)^{-1}}
 \ge c r_i(t/r_i)^{1-\alpha}.
\]
Membership in the cutoff support bounds the left side by $3r_i$, so
$t/r_i$ is bounded. For $t>1$ in $B_1(p)$, the initial harmonic
approximation at scale one gives the same conclusion by comparison
with $R_0(y)$. Packing
therefore gives $|\mathcal A_i(x)|\le N$ for a fixed $N$.
For sufficiently small flatness, a net point with $d(x,y)\le r_i$
satisfies $|P_x-P_y|_{R_i(y)^{-1}}<2r_i$, so that
$\phi_{i,y}(P_x)=1$.  Choose $L>C_0+2$ in
\eqref{re:patch-matrix}.  In particular the mass ball of this
weight-one center contains $B_{r_i}(x)$.

These facts justify the application of
\eqref{re:cumulative-input} to
\begin{equation}
 \label{re:active-sums}
 Z_i(x)=\sum_{y\in\mathcal A_i(x)}D_{i,y},\qquad
 S_i(x)=\sum_{y\in\mathcal A_i(x)}
                     \phi_{i,y}(P_x)D_{i,y}.
\end{equation}
We fix these sums from now on.  They are used only pointwise as
quadratic forms; no spatial differentiability of them is asserted.

For a patch let
\[
 T_{i,y}=\{x\in B_1(p):P_x\in\supp\phi_{i,y}\}.
\]
It is nonempty, since it contains $y$.  All its points lie within
$C_0r_i$ of $y$.  By the positive floor in
\eqref{re:pinching-input}, its infimum of $E_i$ is positive.
Choose one point $z_{i,y}\in T_{i,y}$ satisfying
\begin{equation}
 \label{re:fixed-center}
 E_i(z_{i,y})<2\inf_{x\in T_{i,y}}E_i(x),
 \qquad c_{i,y}=P_{z_{i,y}}.
\end{equation}
Thus a single fixed point satisfies
$E_i(z_{i,y})\le2E_i(x)$ at every active testing point of the patch,
including the support of every derivative of its cutoff.  No continuity
of $E_i$ or measurable selection is needed.  Restricting the minimization
to $T_{i,y}$ also keeps its center in the stated domain at the coarsest
scales.

The following is the covector field used in the iteration:
\begin{equation}
 \label{re:field-definition}
 v_i(P)=\sum_{y\in\mathcal Y_i}
           \phi_{i,y}(P)D_{i,y}(P-c_{i,y}).
\end{equation}
It is a smooth field on all of value space, since the sum is finite.

\begin{lemma}[Field estimates]
 \label{re:field-estimates}
At $P=P_x$, $x\in B_1(p)$, write $r=r_i$, $H=H_i(x)$,
$Q=Q_i(x)$, $Z=Z_i(x)$, $S=S_i(x)$, and $E=E_i(x)$.
There is a linear error $e(\xi)$ such that
\begin{align}
 r^{-1}|v_i|_Q&\le C\sqrt\varepsilon\sqrt E,
       \label{re:field-value}\\
 (\mathrm Dv_i)\xi&=S\xi+e(\xi),
       \label{re:field-linear}\\
 |S\xi|_Q&\le\sqrt\varepsilon\sqrt{Z[\xi]},
       &|e(\xi)|_Q&\le C\sqrt\varepsilon\sqrt E\,|\xi|_H,
       \label{re:field-small-linear}\\
 |\xi^Te(\xi)|&\le C\sqrt{Z[\xi]E}\,|\xi|_H,
       \label{re:field-diagonal}\\
 r|\mathrm D^2v_i(\xi,\xi)|_Q
       &\le C\left(\sqrt{Z[\xi]}\,|\xi|_H
                          +\sqrt E\,|\xi|_H^2\right).
       \label{re:field-second}
\end{align}
There is a fixed $C_v\ge1$ such that, in every order $m\ge0$,
\begin{equation}
 \label{re:field-jets}
 r_i^{m-1}\|\mathrm D^m v_i(P_x)\|_{H_i(x)\to Q_i(x)}
       \le\varepsilon C_v^{m+1}(m!)^s.
\end{equation}
\end{lemma}

\begin{proof}
All sums in this proof are over $\mathcal A_i(x)$.  Put
$w_y=P_x-c_{i,y}$.  Equation~\eqref{re:directional-input} and the
fixed choice~\eqref{re:fixed-center} give
\begin{equation}
 \label{re:center-energy}
 D_{i,y}[w_y]\le C r^2E.
\end{equation}
The same center-distance bounds and harmonic normalization give
$|w_y|_H\le Cr$.  Since $0\le D_{i,y}\le Z\le\varepsilon H$,
conjugation by $H^{-1/2}$ gives
\begin{equation}
 \label{re:positive-matrix-bound}
 |D_{i,y}w|_Q^2\le\varepsilon D_{i,y}[w]
 \quad\text{and}\quad
 |\xi^TD_{i,y}w|^2\le D_{i,y}[\xi]D_{i,y}[w].
\end{equation}
Here the second inequality is Cauchy--Schwarz for a positive
semidefinite form and does not require its invertibility.  Applying
the first inequality to~\eqref{re:center-energy} and using the bounded
number of summands proves~\eqref{re:field-value}.

Differentiation of~\eqref{re:field-definition} gives
\[
 e(\xi)=\sum_y
          (\mathrm D\phi_{i,y})\xi\ D_{i,y}w_y.
\]
By normalized cutoff bounds and matrix comparability,
\begin{equation}
 \label{re:cutoff-derivative}
 |(\mathrm D\phi_{i,y})\xi|\le Cr^{-1}|\xi|_H,
 \qquad
 |\mathrm D^2\phi_{i,y}(\xi,\xi)|\le Cr^{-2}|\xi|_H^2.
\end{equation}
These inequalities and~\eqref{re:positive-matrix-bound} prove the
bound for $e$ in~\eqref{re:field-small-linear}.  Since
$0\le S\le Z\le\varepsilon H$, the same conjugation gives
$SQS\le\varepsilon S\le\varepsilon Z$, proving its bound for $S$.
For the sharper diagonal bound, use the second inequality in
\eqref{re:positive-matrix-bound} before summing:
\[
 |\xi^Te(\xi)|
 \le C\sqrt E\,|\xi|_H\sum_y\sqrt{D_{i,y}[\xi]}
 \le C\sqrt{Z[\xi]E}\,|\xi|_H.
\]
Finally,
\[
 \mathrm D^2v_i(\xi,\xi)
 =\sum_y\left(
   \mathrm D^2\phi_{i,y}(\xi,\xi)D_{i,y}w_y
       +2(\mathrm D\phi_{i,y})\xi\ D_{i,y}\xi\right).
\]
Equations~\eqref{re:positive-matrix-bound}--\eqref{re:cutoff-derivative}
bound the two sums by
$Cr^{-1}\sqrt\varepsilon\sqrt E\,|\xi|_H^2$ and
$Cr^{-1}\sqrt\varepsilon\sqrt{Z[\xi]}\,|\xi|_H$, respectively.
Dropping $\sqrt\varepsilon\le1$ proves~\eqref{re:field-second}.

For the isotropic estimates, $D_{i,y}\le\varepsilon H$ also gives
$\|D_{i,y}\|_{H\to Q}\le\varepsilon$ and
$|D_{i,y}w_y|_Q\le C\varepsilon r$.  In order $m\ge1$, the
product in~\eqref{re:field-definition} has only the term with $m$
derivatives on the cutoff and the $m$ terms with one derivative on
the affine factor.  By~\eqref{re:bump-jets}, the frozen source-norm
comparison, and bounded overlap, their total norm is at most
\[
 C\varepsilon r^{1-m}
 \left(C_1^{m+1}(m!)^s
                  +m C_1^m((m-1)!)^s\right).
\]
Increase $C_v$ once to obtain~\eqref{re:field-jets}.  The preceding
value bound with $|w_y|_H\le Cr$ proves its $m=0$ case.
\end{proof}

\subsection{The iteration and its derivative bounds}

On the common open set $\Omega$ define, as long as the Jacobians are
invertible there,
\begin{equation}
 \label{re:iteration}
 h_0(P)=P,
 \qquad J_i(P)=\mathrm Dh_i(P),
 \qquad h_{i+1}(P)=h_i(P)+J_i(P)^{-T}v_i(P).
\end{equation}
The inverse transpose is chosen to turn the leading derivative of
$v_i$ directly into a change of metric.  To see both the gain and the
error, put $W=J^{-T}v$ for a smooth map $h$ with $J=\mathrm Dh$.
Differentiating $J^TJ^{-T}=\Id$ gives the exact identity
\begin{equation}
 \label{re:metric-step}
 \begin{aligned}
 |(J+\mathrm DW)\xi|^2-|J\xi|^2
   ={}&2\xi^T\mathrm Dv\,\xi
       -2\langle\mathrm D^2h(\xi,\xi),J^{-T}v\rangle\\
     &+|\mathrm DW\,\xi|^2.
 \end{aligned}
\end{equation}
By the field estimates, $\mathrm Dv_i\,\xi=S_i\xi+e(\xi)$.
The first term therefore supplies the desired positive increment
$2S_i[\xi]$, up to the directional error.  The middle term explains
why we will estimate the diagonal derivative
$\mathrm D^2h_i(\xi,\xi)$ by a scale recurrence.  Bounds on all
higher derivatives serve a different purpose: they keep the smooth
update controlled despite its loss of one derivative.

We shall prove simultaneously a metric bound and all higher-derivative
bounds.  Choose $A\ge1$, subsequently fixed in terms of $C_*$, and
fix $\sigma\ge\max\{1,s\}$.  The desired bounds are
\begin{align}
 A^{-1}H_i(x)&\le J_i(P_x)^TJ_i(P_x)\le A H_i(x),
      \label{re:metric-bootstrap}\\
 r_i^{m-1}\|\mathrm D^m h_i(P_x)\|_{H_i(x)\to\R^n}
       &\le B^{m-1}(m!)^{1+\sigma}\qquad(m\ge2),
      \label{re:jet-bootstrap}
\end{align}
where $B$ will depend on the fixed data and $A,\sigma$, but not on
$\varepsilon,i,x,m$.

The next lemma shows that matrix inversion preserves the exponential
base $B$ in these derivative bounds.  A new factor growing exponentially
with the order at each step would prevent the simultaneous induction.

\begin{lemma}[Factorial convolution and inverse coefficients]
 \label{re:inverse-majorant}
Fix $\sigma\ge1$ and put $L_0=1$ and
\[
 L_m=(m+1)((m+1)!)^\sigma\qquad(m\ge1).
\]
Then
\begin{equation}
 \label{re:convolution-decay}
 \delta_m:=\sum_{l=1}^{m-1}\frac{L_lL_{m-l}}{L_m}
       \longrightarrow0,
 \qquad
 \sum_{l=0}^m L_lL_{m-l}\le C_\sigma L_m.
\end{equation}
Suppose a smooth matrix-valued function $J$ is invertible at a point
and, in Euclidean operator norms at that point,
\[
 \|J^{-1}\|\le\sqrt A,\qquad
 \frac{\|\mathrm D^mJ\|}{m!}\le B^mL_m\quad(m\ge1),
 \qquad A,B\ge1.
\]
There is a constant $C_{A,\sigma}$, independent of $m$ and $B$, such that
\begin{equation}
 \label{re:inverse-jets}
 \frac{\|\mathrm D^m(J^{-1})\|}{m!}
       \le C_{A,\sigma}B^mL_m\qquad(m\ge0).
\end{equation}
The assertion also holds for the inverse transpose.
\end{lemma}

\begin{proof}
For $1\le l\le m-1$ the summand in the first expression of
\eqref{re:convolution-decay} is
\[
 \frac{(l+1)(m-l+1)}{m+1}
 \left(\frac{(l+1)!(m-l+1)!}{(m+1)!}\right)^\sigma.
\]
The two terms with $l=1,m-1$ are $O_\sigma(m^{-\sigma})$;
the terms with $l=2,m-2$ are $O_\sigma(m^{-2\sigma})$.
For $3\le l\le m-3$, the factorial ratio is at most
$24/((m+1)m(m-1))$: its maximum in this interval occurs at an
endpoint, as is seen by taking the ratio of consecutive terms.
The first factor is at most $m+1$, and there are at most $m$ such
terms.  Their sum is therefore $O_\sigma(m^{2-3\sigma})$,
which is $O_\sigma(m^{-\sigma})$ because $\sigma\ge1$.
This proves $\delta_m\to0$; the finitely many small values of $m$
do not matter.  Including the terms $l=0,m$ proves the second
assertion of~\eqref{re:convolution-decay}.

Let $K=J^{-1}$.  The multilinear product formula, followed by the
operator norm inequality and division by $m!$, gives
\[
 \frac{\|\mathrm D^m K\|}{m!}
 \le\sqrt A\sum_{l=1}^m
        \frac{\|\mathrm D^l J\|}{l!}
        \frac{\|\mathrm D^{m-l}K\|}{(m-l)!}.
\]
In particular the binomial coefficients in the product formula have
been removed exactly.  After dividing out $B^m$, universal majorants
$a_m$ for the positive-order coefficients satisfy
\begin{equation}
 \label{re:inverse-recursion}
 a_m\le A L_m+\sqrt A
                    \sum_{l=1}^{m-1}L_l a_{m-l}.
\end{equation}
Choose $M=M(A,\sigma)$ so that $\sqrt A\,\delta_m\le1/2$ for
$m\ge M$.  The recursion for the finitely many orders $m<M$
gives finite universal bounds depending only on $A$ and $\sigma$.
Choose $C_{A,\sigma}$ to dominate $\sqrt A$, $2A$, and all these
bounds after division by $L_m$.  If $a_q\le C_{A,\sigma}L_q$
for $q<m$, then~\eqref{re:inverse-recursion} gives, for $m\ge M$,
\[
 \frac{a_m}{L_m}
       \le A+\sqrt A\,C_{A,\sigma}\delta_m
       \le A+\frac12 C_{A,\sigma}
       \le C_{A,\sigma}.
\]
This proves~\eqref{re:inverse-jets} with a single constant for all
orders; no factor $C_{A,\sigma}^m$ has been introduced.  Transposition
preserves Euclidean operator norms and the same derivative estimates.
\end{proof}

\begin{lemma}[One step for all jets]
 \label{re:jet-step}
Choose $B\ge C_v$ sufficiently large and then $\varepsilon$ sufficiently
small, depending on $A,\sigma,B$ and the fixed data.  If
\eqref{re:metric-bootstrap} and~\eqref{re:jet-bootstrap} hold at level
$i$, the update~\eqref{re:iteration} satisfies
\eqref{re:jet-bootstrap} at level $i+1$.  This assertion does not
require the new-level metric bound.
\end{lemma}

\begin{proof}
Freeze $x$ and abbreviate $r=r_i$, $H=H_i(x)$, $Q=H^{-1}$.
The linear source and target normalizations are
\[
 \begin{aligned}
 P&=P_x+rH^{-1/2}X,\\
 \widetilde h(X)
   &=r^{-1}\bigl(h_i(P_x+rH^{-1/2}X)-h_i(P_x)\bigr),\\
 \widetilde v(X)&=r^{-1}H^{-1/2}v_i(P_x+rH^{-1/2}X).
 \end{aligned}
\]
At $X=0$, $\widetilde J=\mathrm D\widetilde h=J_iH^{-1/2}$
has singular values between $A^{-1/2}$ and $A^{1/2}$, and
\[
 \widetilde J^{-T}\widetilde v
           =r^{-1}J_i^{-T}v_i.
\]
This identity checks the covector normalization in the update.  All matrices in the change of variables
are constant matrices.

For every $m\ge1$,~\eqref{re:jet-bootstrap} gives
\[
 \frac{\|\mathrm D^m\widetilde J(0)\|}{m!}
 \le\frac{B^m((m+1)!)^{1+\sigma}}{m!}
 =B^m L_m.
\]
Lemma~\ref{re:inverse-majorant} therefore bounds the coefficients
of $\widetilde J^{-T}$ by $C_{A,\sigma}B^mL_m$.
On the other hand,~\eqref{re:field-jets} gives
\[
 \frac{\|\mathrm D^q\widetilde v(0)\|}{q!}
       \le\varepsilon C_v^{q+1}(q!)^{s-1}
       \le\varepsilon C_v B^qL_q\qquad(q\ge0).
\]
The product rule divided by $m!$ and
\eqref{re:convolution-decay} now imply
\begin{equation}
 \label{re:increment-jets}
 \frac{\|\mathrm D^m(\widetilde J^{-T}\widetilde v)(0)\|}{m!}
       \le\varepsilon C_A B^m L_m.
\end{equation}

To pass to next-scale jet units, note that
$H_{i+1}(x)\ge H_i(x)$.  Thus the norm of any multilinear map with
source norm $H_{i+1}(x)$ is at most its norm with source norm $H_i(x)$.
The change from $r_i$ to $r_{i+1}=r_i/2$ supplies the additional factor
$2^{1-m}$.  Relative to the allowance
$B^{m-1}(m!)^{1+\sigma}$ at order $m\ge2$, the old contribution is
at most $2^{1-m}\le1/2$, while~\eqref{re:increment-jets} contributes
at most
\[
 \varepsilon C_A B\,2^{1-m}(m+1)^{1+\sigma}.
\]
The supremum of $2^{1-m}(m+1)^{1+\sigma}$ over $m\ge2$ is finite.
One sufficiently small choice of $\varepsilon$ therefore makes the
new contribution at most one half in every order simultaneously.
This proves the lemma.
\end{proof}

\subsection{Diagonal second derivatives and the metric bootstrap}

\begin{proposition}[Uniform metric and jet bounds]
 \label{re:metric-iteration}
There are fixed $A,B$ and a sufficiently small flatness threshold
such that the maps~\eqref{re:iteration} are smooth on all of $\Omega$
for every finite $i$, their Jacobians are invertible there, and
\eqref{re:metric-bootstrap}--\eqref{re:jet-bootstrap} hold for all
$x\in B_1(p)$ and every $i$.
\end{proposition}

\begin{proof}
We give the estimates under the induction hypothesis through level
$k-1$ and then close that induction.  Fix $x$ and a nonzero source
vector $\xi$.  Every expression in this proof is evaluated at $P_x$.
Set
\[
 q_i=|\xi|_{H_i(x)},\qquad
 b_i=\frac{r_i|\mathrm D^2h_i(\xi,\xi)|}{q_i},\qquad
 K_i=J_i^{-T},\qquad W_i=K_iv_i.
\]
The metric bound gives $|K_iw|\le\sqrt A|w|_{Q_i(x)}$ for
covectors $w$, since
$(J_i^TJ_i)^{-1}\le A Q_i(x)$.
The jet bounds through order three, the ordinary inverse derivative
formula, and the duality of $H_i$ and $Q_i$ give
\begin{align}
 |(\mathrm D_\xi K_i)w|
     &\le C_A\frac{q_i}{r_i}|w|_{Q_i(x)},
       \label{re:first-inverse-cost}\\
 |(\mathrm D_\xi^2 K_i)w|
     &\le C_A\frac{q_i^2}{r_i^2}|w|_{Q_i(x)}.
       \label{re:second-inverse-cost}
\end{align}
For completeness, after the frozen normalization in the preceding
lemma, the first formula is
$\mathrm D_\xi K_i=-K_i(\mathrm D_\xi J_i)^TK_i$.
The second derivative is a sum of two products containing two first
derivatives of $J_i$ and one product containing
$\mathrm D_\xi^2J_i$.  Bounds for $\mathrm D^2h_i$ and
$\mathrm D^3h_i$, with the bounded normalized inverse, give exactly
\eqref{re:first-inverse-cost}--\eqref{re:second-inverse-cost} after
undoing the normalization.  Thus these estimates are also valid for
arbitrarily anisotropic $H_i$.

Twice differentiating $W_i$ gives all three types of terms:
\begin{equation}
 \label{re:second-increment-terms}
 \mathrm D^2W_i(\xi,\xi)
   =K_i\mathrm D^2v_i(\xi,\xi)
     +2(\mathrm D_\xi K_i)(\mathrm Dv_i\xi)
     +(\mathrm D_\xi^2 K_i)v_i.
\end{equation}
By Lemma~\ref{re:field-estimates} and
\eqref{re:first-inverse-cost}--\eqref{re:second-inverse-cost},
multiplication by $r_i/q_i$ bounds each of these terms by a constant
times
\[
 f_i:=\sqrt{Z_i(x)[\xi]}+\sqrt{E_i(x)}\,q_i.
\]
More explicitly, the first term uses~\eqref{re:field-second}; the
second uses $|\mathrm Dv_i\xi|_{Q_i}\le
C\sqrt\varepsilon f_i$; and the third uses
$|v_i|_{Q_i}\le C\sqrt\varepsilon r_i\sqrt{E_i(x)}$.
Because $q_{i+1}\ge q_i$ and $r_{i+1}=r_i/2$, the old second
derivative contracts by at least a factor of one half.  Hence
\begin{equation}
 \label{re:diagonal-recurrence}
 b_{i+1}\le\tfrac12 b_i+C_A f_i,
 \qquad b_0=0.
\end{equation}
Convolution with the sequence $1,1/2,1/4,\ldots$, whose sum is two,
is bounded on $\ell^2$.  Thus, for every finite $k$ covered by the
induction hypothesis,
\begin{align}
 \sum_{i<k}b_i^2
 &\le C_A\sum_{i<k}\left(Z_i(x)[\xi]+E_i(x)q_i^2\right)
       \notag\\
 &\le C_A H_k(x)[\xi].
       \label{re:diagonal-square-sum}
\end{align}
The last line uses~\eqref{re:cumulative-input}, monotonicity
$H_i\le H_k$, and~\eqref{re:pinching-input}.  In particular it
does not require that the matrices $H_i$ remain bounded as
$i\to\infty$.

The first derivative of the increment satisfies
\[
 \Delta_iJ\,\xi:=(J_{i+1}-J_i)\xi
       =K_i(\mathrm Dv_i\xi)+(\mathrm D_\xi K_i)v_i.
\]
The same estimates give the pointwise square bound
\begin{equation}
 \label{re:squared-increment}
 |\Delta_iJ\,\xi|^2
 \le C_A\varepsilon
           \left(Z_i(x)[\xi]+E_i(x)q_i^2\right).
\end{equation}
Consequently their sum for $i<k$ is at most
$C_A\varepsilon H_k(x)[\xi]$.

It remains to control the linear part of the increment of the squared
length.  The undifferentiated inverse gives
\[
 \langle J_i\xi,K_i\mathrm Dv_i\xi\rangle
       =\xi^T\mathrm Dv_i\xi
       =S_i(x)[\xi]+\xi^Te_i(\xi).
\]
For the differentiated inverse, direct matrix multiplication gives
\begin{align}
 \langle J_i\xi,(\mathrm D_\xi K_i)v_i\rangle
 &=-\xi^T(\mathrm D_\xi J_i)^TJ_i^{-T}v_i
       \notag\\
 &=-\langle\mathrm D^2h_i(\xi,\xi),K_iv_i\rangle_{\R^n}.
       \label{re:inverse-transpose-contraction}
\end{align}
The inner product on the right is the ordinary Euclidean target
inner product.  Its absolute value is bounded by
$C_A b_iq_i\sqrt{E_i(x)}$, using~\eqref{re:field-value} and
dropping the additional factor $\sqrt\varepsilon\le1$.

Now expand $|J_{i+1}\xi|^2-|J_i\xi|^2$, sum over $i<k$,
and use~\eqref{re:field-diagonal},
\eqref{re:squared-increment}, and
\eqref{re:inverse-transpose-contraction}.  The two linear error
sums satisfy, respectively,
\begin{align*}
 \sum_{i<k}|\xi^Te_i(\xi)|
 &\le C\left(\sum_{i<k}Z_i(x)[\xi]\right)^{1/2}
          \left(\sum_{i<k}E_i(x)q_i^2\right)^{1/2}
 \le C\sqrt\varepsilon H_k(x)[\xi],\\
 \sum_{i<k}b_iq_i\sqrt{E_i(x)}
 &\le\left(\sum_{i<k}b_i^2\right)^{1/2}
          \left(\sum_{i<k}E_i(x)q_i^2\right)^{1/2}
 \le C_A\sqrt\varepsilon H_k(x)[\xi].
\end{align*}
We have proved the quantitative cumulative error estimate
\begin{equation}
 \label{re:metric-error}
 \left|\,|J_k\xi|^2-|\xi|^2
                       -2\sum_{i<k}S_i(x)[\xi]\,\right|
       \le C_A(\sqrt\varepsilon+\varepsilon)H_k(x)[\xi].
\end{equation}
Only the sum of $E_i$ has been used, never the sum of its square
roots.

We now fix the constants and complete the simultaneous induction.
By~\eqref{re:cumulative-input} and $0\le S_i\le Z_i$,
\begin{equation}
 \label{re:leading-metric-comparison}
 C_*^{-1}H_k
       \le\Id+2\sum_{i<k}S_i
       \le2C_*H_k.
\end{equation}
Choose $A\ge4C_*$, increasing it if necessary for the initial
comparability of $H_0$ and $\Id$.  Fix $\sigma,B$ as in
Lemma~\ref{re:jet-step}.  Only now choose the data sufficiently flat
that its smallness requirement holds and
\[
 C_A(\sqrt\varepsilon+\varepsilon)<\frac1{2C_*}.
\]
Then~\eqref{re:metric-error} and
\eqref{re:leading-metric-comparison} imply
\eqref{re:metric-bootstrap} with slack at level $k$.

Initially $h_0$ is affine, so all its positive-order derivatives
beyond the first vanish, and both bootstrap statements hold.
Suppose they hold through level $i$.  The pointwise inverse of $J_i$
exists everywhere on $\Omega$ and is smooth there; hence the update
defines a smooth $h_{i+1}$ on this same open set.  The proof of
\eqref{re:metric-error} at $k=i+1$ uses inverse and jet estimates
only through level $i$: the error sum uses $b_j$ only for $j\le i$,
whose recurrence involves still earlier increments.  It therefore
proves the metric bound and invertibility at level $i+1$ without
assuming them.  Lemma~\ref{re:jet-step} proves all jet bounds at
level $i+1$ using the old-level data.  This is an induction on the
finite scale index with an all-orders assertion at each step, and
completes the proof.
\end{proof}

\begin{remark}
The construction neither differentiates the fields $H_i$ nor assumes
that they have a finite limit.  It also does not require the maps
$h_i$ themselves to be globally injective.  Their ordinary Jacobians
are invertible on the fixed open value-space domain, which is exactly
what the update needs.
\end{remark}

\subsection{Passage to an open bi-Lipschitz chart}

\begin{proposition}[The limiting chart]
 \label{re:limit-chart}
For sufficiently flat data in the preceding construction, there is
$\rho>0$ such that the uniform limit
\[
 F(x)=\lim_{i\to\infty}h_i(u(x)),\qquad x\in B_1(p),
\]
restricts to a bi-Lipschitz homeomorphism from $B_\rho(p)$ onto an
open subset of $\R^n$.  Its distance comparisons use the ambient
distance $d$ restricted to $B_\rho(p)$.
\end{proposition}

\begin{proof}
From~\eqref{re:field-value} and
\eqref{re:metric-bootstrap},
\[
 |h_{i+1}(P_x)-h_i(P_x)|
       \le C_A\sqrt\varepsilon r_i\sqrt{E_i(x)}
       \le C_A\varepsilon r_i.
\]
The geometric series gives uniform convergence on $B_1(p)$ and a
tail estimate
\begin{equation}
 \label{re:uniform-tail}
 \sup_{x\in B_1(p)}|F(x)-h_i(P_x)|\le\tau r_i,
 \qquad \tau=C_A\varepsilon\longrightarrow0.
\end{equation}
Each approximating map is continuous, so $F$ is continuous.

Let $K_2=B(2!)^{1+\sigma}$, the fixed second-derivative constant in
\eqref{re:jet-bootstrap}.  By~\eqref{re:ellipsoid-input}, along
any value-space segment contained in the central ellipsoid at $x$,
with preimages in $B_{r_i}(x)\subset B_1(p)$,
\begin{equation}
 \label{re:taylor-constant}
 |\mathrm D^2h_i(P_y)(\zeta,\zeta)|
       \le\frac{K_2C_{\mathrm{cmp}}}{r_i}
                         |\zeta|_{H_i(x)}^2.
\end{equation}
Put $K_T=K_2C_{\mathrm{cmp}}/2$.  Choose a fixed $D\ge1$ so large
that
\begin{equation}
 \label{re:choice-D}
 D>4/a,\qquad 4K_T/D\le1/(8\sqrt A).
\end{equation}
These choices depend only on constants already fixed in the
construction.

Take a further sufficiently late rescaling, without changing those
constants, so that the normalized harmonic approximation at the
fixed scale ratio $D$ gives
\begin{equation}
 \label{re:pair-normalization}
 \tfrac12 d(x,z)\le|P_z-P_x|_{H_i(x)}\le2d(x,z)
 \quad\text{whenever }Dr_i^{-1}d(x,z)\in[1/2,1]
\end{equation}
on the interior balls in use.  To justify this precise consequence,
normalized harmonic approximation controls distances in value space
up to $o(1)r_i$ on the fixed ball; since the displayed distances are
at least $r_i/(2D)$, its relative error tends to zero after $D$ has
been fixed.  The relative comparison of $Q_i$ with the normalization
matrices preserves these bounds.  Also require
\begin{equation}
 \label{re:choice-tail}
 4D\tau\le1/(8\sqrt A).
\end{equation}

Set, for example, $\rho=(16D+4)^{-1}$.  If $x,z\in B_\rho(p)$
are distinct and $t=d(x,z)$, then $t\le2\rho$.  Choose a dyadic
scale with
\[
 Dt\le r_i<2Dt.
\]
It has $r_i<4D\rho<1$, and
$B_{r_i}(x)\subset B_1(p)$.  Thus every pair in the final ball,
not merely sufficiently generic pairs, is covered by an available
scale.  Write $\Delta=P_z-P_x$.  Equation~\eqref{re:pair-normalization}
gives $t/2\le|\Delta|_{H_i(x)}\le2t$.
By~\eqref{re:choice-D}, $|\Delta|_{H_i(x)}<a r_i$.
The entire straight segment
\[
 P_x+s\Delta,\qquad 0\le s\le1,
\]
therefore lies in the central ellipsoid in
\eqref{re:ellipsoid-input}, and hence in
$u(B_{r_i}(x))\subset\Omega$.  In particular every point where a
derivative is integrated has an interior preimage and obeys the
matrix comparison used in~\eqref{re:taylor-constant}.

Taylor's formula along this segment yields
\[
 h_i(P_z)-h_i(P_x)=J_i(P_x)\Delta+\mathcal R,
 \qquad
 |\mathcal R|\le\frac{K_T}{r_i}|\Delta|_{H_i(x)}^2
       \le\frac{4K_T}{D}t\le\frac{t}{8\sqrt A}.
\]
The metric bound gives
\[
 \frac{t}{2\sqrt A}\le|J_i(P_x)\Delta|\le2\sqrt A\,t.
\]
Finally the two tails in~\eqref{re:uniform-tail} have total size at
most $2\tau r_i\le4D\tau t\le t/(8\sqrt A)$.  Combining the
last three inequalities proves, for every $x,z\in B_\rho(p)$,
\begin{equation}
 \label{re:final-bilipschitz}
 \frac{1}{4\sqrt A}d(x,z)
       \le|F(x)-F(z)|\le4\sqrt A\,d(x,z).
\end{equation}
The case $x=z$ is immediate.  In particular $F$ is injective there
and has a continuous inverse onto its image.

The set $u(B_\rho(p))$ is open, by the previously established open
coordinates.  The continuous injective map
\[
 F\circ u^{-1}:u(B_\rho(p))\longrightarrow\R^n
\]
therefore has open image by invariance of domain
\cite[Theorem~2B.3]{Hatcher02}. This proves the
proposition.  All scale choices used $t=d(x,z)$ in the ambient space;
no intrinsic path distance on the chosen neighborhood was substituted.
\end{proof}

\begin{proof}[Proof of Theorem~\ref{thm:main}]
Let $p$ be the prescribed $n$-regular point of the noncollapsed
$\RCD(K,n)$ space.  The preceding sections provide the geometric,
harmonic, heat, and directional-pinching conclusions on sufficiently
late rescalings at $p$.  First fix the finitely many patch enlargement
constants using only $n$. Choose the cutoff with $s=2$ and fix
$\sigma=2$. The cumulative constant $C_*$, the bootstrap constants
$A,B$, and then $D$ in \eqref{re:choice-D} consequently depend only
on $n$. Choose a rescaling sufficiently late that all
their flatness requirements, the simultaneous induction, and
\eqref{re:choice-tail} hold.  Proposition~\ref{re:limit-chart} gives
an open ball about $p$ and a chart with the required two-sided
distance bound on that rescaled space.  If its metric is $d/r_*$,
multiply the coordinate values of that chart by $r_*$ when returning
to the original metric.  Both sides of
\eqref{re:final-bilipschitz} then scale by the same factor, while
the domain remains an open neighborhood and the image remains an
open subset of $\R^n$. Thus Theorem~\ref{thm:main} holds with
$L_n=4\sqrt{A(n)}$. The required rescaling, and hence the radius
of the neighborhood in the original metric, may depend on $X$ and $p$.
\end{proof}
